\section{Conclusions and Future Work}
\label{sec:conclusion}

\subsection{Answer to the central research question}

This project examined whether published diagnostic and prognostic workflows for
liver metastasis in pancreatic neuroendocrine tumours could be reconstructed
using the available Surveillance, Epidemiology, and End Results (SEER) data. It
also asked whether Tabular Prior-Data Fitted Network (TabPFN) representations,
synthetic data, multilayer perceptrons (MLPs) and other tabular foundation models
could improve prediction.

The reference workflows could be reconstructed and produced strong diagnostic
discrimination and meaningful prognostic risk ordering in the available SEER
cohorts. Exact replication of the published results was not possible because
the original SEER release and external cohorts were no longer available. The
independent University Hospital Erlangen (UK Erlangen) cohort provided an
external evaluation of the reproduced reference models. Their predictive
ranking remained encouraging, but the cohort of 30 diagnostic and 11
prognostic patients was too small to estimate performance precisely or support
claims of clinical utility.

The additional modelling experiments produced a more selective answer. Three
tabular foundation models achieved small diagnostic improvements over logistic
regression when they were evaluated on the same SEER test patients. The tested
MLP did not provide a general improvement, and none of the foundation models or
the MLP improved prognostic discrimination reliably. Synthetic augmentation
benefited diagnostic XGBoost, but the improvement did not extend consistently
to the other diagnostic models or to prognosis.

The central conclusion is therefore that foundation models and synthetic data
can provide useful gains in specific settings, but their benefit is not
automatic. Performance depends on the prediction task, the model, the amount
and quality of real data, the evaluation design and the clinical relationships
preserved by the transformation. Generated records and more complex models did
not replace the need for appropriate survival methods or independent patients.

\subsection{Main scientific conclusions}

\paragraph{Reconstruction of the reference study.}
The eligibility criteria and preprocessing procedure based on the reference
study produced a diagnostic cohort of 2,672 patients and a prognostic cohort of
462 patients with liver metastasis. The reconstructed diagnostic and prognostic
analyses recovered the main predictive patterns reported in the publication.
Numerical differences remain expected because the available SEER release,
patients and some variable definitions differ from those of the original study.

\paragraph{External evaluation of the reproduced models.}
In the common set of 29 UK Erlangen diagnostic patients, the primary
five-variable nomogram achieved an area under the receiver operating
characteristic curve (ROC AUC) of 0.934, with a 95\% confidence interval from
0.809 to 1.000. In all 11 prognostic patients, the primary five-variable Cox
nomogram achieved a Harrell concordance index of 0.833, with a confidence
interval from 0.625 to 1.000. These results show that the unchanged SEER-trained
models retained predictive ranking in the available hospital cohort. The wide
intervals, small calibration groups and survival groups of only four and seven
patients prevent conclusions about stable performance or clinical benefit.

\paragraph{Diagnostic model comparison.}
TabPFN-3, TabICLv2 and TabDPT achieved ROC AUCs of 0.902--0.903 in the repeated
SEER evaluation. Their improvements over logistic regression ranged from
0.0074 to 0.0088 when the models were compared on the same patients. These are
small internal gains rather than evidence of clinically important superiority.
The clinical MLP performed competitively but did not improve consistently over
the simpler comparison models. Adding a generic nonlinear prediction model was
therefore not sufficient to reproduce the strongest diagnostic results.

\paragraph{Prognostic model comparison.}
No foundation model or MLP showed a reliable concordance improvement in
prognosis. Random survival forest achieved the lowest integrated Brier score
(IBS) of 0.153 and improved IBS by 0.0071 relative to Cox elastic net when both
models were evaluated on the same patients. This modest probability-error
advantage favoured an established survival method. More generally, the
prognostic results show that improvements in binary diagnosis should not be
assumed to transfer to a smaller, censored survival task.

\paragraph{Representations and synthetic data.}
The TabPFN embeddings retained predictive information and much of the local
patient-neighbourhood structure. However, they were created from the complete
cohorts using outcome information, so their downstream results remain
exploratory. The complete synthetic datasets reproduced many individual
variable distributions and allowed the diagnostic and prognostic pipelines to
run, but grade 4 was not generated, survival distributions shifted and several
relationships between clinical variables were not fully preserved. These
datasets are useful for studying representation and fidelity; they are not
independent clinical evidence.

\paragraph{Synthetic augmentation.}
For each evaluation partition, synthetic records were generated using only the
real training patients, and performance was measured only in held-out real
patients. At the largest augmentation size of 15,000 records, diagnostic
XGBoost improved ROC AUC by 0.0126 and reduced Brier score by 0.0060 compared
with training on real data alone. Other diagnostic models did not improve
consistently, and prognostic augmentation provided no benefit. Some strong
models performed worse as synthetic records formed a larger share of the
training data. Synthetic sample size should therefore not be interpreted as an
equivalent increase in the number of observed patients.

\subsection{Contributions of the project}

Within the limitations described in Section~\ref{sec:limitations}, the project
makes five main contributions:

\begin{enumerate}[label=\textbf{C\arabic*.},ref=C\arabic*]
\item It constructs clearly defined diagnostic and prognostic cohorts from the
SEER registry using eligibility, coding and preprocessing decisions based on
the reference study.

\item It reconstructs the reference diagnostic and prognostic workflows on the
closest available SEER population and evaluates the resulting models without
modification in an independent UK Erlangen cohort.

\item It creates and evaluates TabPFN embeddings and complete synthetic
datasets for predictive information, structural similarity, category support
and preservation of clinically relevant relationships.

\item It compares conventional models, tabular foundation models and MLPs on
the same real-patient evaluation partitions, enabling direct comparison between
model families.

\item It tests synthetic augmentation by generating records only from the real
training patients and evaluating models in held-out real patients, separating
practical utility from descriptive similarity.
\end{enumerate}

These contributions do not constitute a deployed clinical system, a causal
analysis of treatment effects or a formal privacy guarantee.

\subsection{Priorities for future research}

Future work should first strengthen the clinical evidence.

\begin{enumerate}[label=\textbf{P\arabic*.},ref=P\arabic*]
\item \textbf{Larger external validation.} Evaluate the most promising
approaches in larger multicentre and temporally distinct cohorts, with outcomes
and predictor timing specified before the models are applied.

\item \textbf{Training-only transformations and complete uncertainty.} Learn
embeddings, select variables, train synthetic-data generators and calibrate
predictions using training patients only. Repeat the complete development
process to measure both construction and evaluation uncertainty.

\item \textbf{Clinically constrained synthetic data.} Preserve rare categories
such as grade 4, enforce agreement between linked fields and use methods
designed for censored outcomes. Evaluate utility, fidelity and privacy risk
together, preferably in independent data.

\item \textbf{Survival-specific foundation models.} Develop and test models
that account for censoring and produce valid survival probabilities across time.
Use larger cohorts with more observed deaths and retain strong conventional
survival methods as comparisons.

\item \textbf{Prospective clinical evaluation.} After adequate external
validation, define thresholds in advance, examine calibration and subgroup
performance, and test whether model use improves clinical decisions or
patient-relevant outcomes.
\end{enumerate}

\subsection{Final conclusion}

Modern tabular foundation models and synthetic data can support rare-tumour
registry research, but greater complexity does not guarantee clinical value.
The strongest findings are small internal diagnostic gains, a selective
augmentation benefit for XGBoost and encouraging but imprecise external results
for the reproduced reference models. No general prognostic, MLP or augmentation
advantage was demonstrated. Larger independent clinical cohorts are now
essential; until then, the models and generated records should remain research
tools rather than validated support for clinical decisions.
